\section{Signal Processing Fundamentals}
\label{app:foundations}

This section reviews the theoretical basis for using complex numbers to represent signals and Jones matrices to represent linear transformations of the electric field.

\subsection{The Analytic Signal}
\label{app:analytic}

The voltage signal from each receptor is a real-valued function of
time, or process, that may be represented by its associated analytic
signal.
%
The analytic signal, also known as Gabor's complex signal, is a
complex-valued representation of a real-valued process that provides
its instantaneous amplitude and phase.  
%
The analytic signal associated with a process, $x(t)$, is defined by
\begin{equation}
  z(t) = x(t) + \Ci\hat x(t),
\end{equation}
%
where $\hat x(t)$ is the Hilbert transform of $x(t)$ \citep{pap65}.
%
Following a Fourier transform, defined by
\begin{equation}
  X(\nu) = \int_\infty^\infty x(t) \exp(- \Ci 2\pi \nu t) dt,
  \label{eqn:Fourier_transform}
\end{equation}
%
the Hilbert transform is equivalent to
\begin{equation}
\hat{X}(\nu)=H(\nu)X(\nu),    
\end{equation}
where
%
\begin{equation} H(\nu) = \left\{ \begin{array}{cc}
                -\Ci  & \nu > 0 \\
                \Ci & \nu < 0
                \end{array} \right.
\end{equation}
is known as the quadrature filter \citep{pap65}.  
%
The quadrature filter can be understood as a
90\degr\ phase shifter by noting that 
the Hilbert transform of $\cos(\nu_0t)$ is equal to
$\sin(\nu_0t)$ (see \Tab{fourier}).  
%
Using the quadrature filter, it can also be
shown that the Fourier transform of the analytic signal, $Z(\nu)$, is
equal to zero for $\nu < 0$.
\begin{equation}
\begin{split}
 Z(\nu) & =  X(\nu) + \Ci\hat X(\nu) \\
	& =  X(\nu) + \Ci H(\nu)X(\nu) \\
	& =  \left \{ \begin{array}{cc} 
	2X(\nu)  & \nu > 0 \\
	0 & \nu < 0 
	\end{array} \right.
\end{split}
\end{equation}
Conversely, the analytic signal associated with $x(t)$ is
produced
by suppressing the negative frequencies in $X(\nu)$.

As it is derived from the real-valued process, the analytic signal
does not contain any additional information.  However, the analytic
signals associated with two orthogonal senses of polarization, $e_0(t)$ and
$e_1(t)$, facilitate calculation of the coherency matrix.  

\begin{table}
\caption
{Useful Fourier Transform pairs. The left column lists functions of
time.  In the right column, the corresponding Fourier transform is
given as a function of oscillation frequency, $\nu$. 
% The filters, $H(\nu)$ and $\Pi(\nu)$, are known as the quadrature and rectangle functions, respectively.
}
\label{tab:fourier}
\begin{center}
\begin{tabular}{c|c}
\hline
\hline
$x(t)$ & $X(\nu)$ \\
\hline \\ [-3mm]
$\cos(2\pi\nu_0t)$ & $(\delta(\nu+\nu_0)+\delta(\nu-\nu_0)) / 2$ \\ [3mm]
$\sin(2\pi\nu_0t)$ & $\Ci (\delta(\nu+\nu_0)-\delta(\nu-\nu_0)) / 2$ \\ [2mm]
%
\hline
\multicolumn{2}{c}{Quadrature Filter} \\
\hline \\ [-3mm]
%
$h(t)=(\pi t)^{-1}$ & 
$H(\nu) =\left\{ \begin{array}{cc}
            -\Ci & \nu > 0 \\
            \Ci  & \nu < 0 \end{array} \right.$ \\ [4mm]
%
\hline
\multicolumn{2}{c}{Lowpass Filter} \\
\hline \\ [-3mm]
%
$\pi(t)=\sinc(\pi\bw t)$ & 
$\Pi(\nu/\bw) = \left\{ \begin{array}{cc}
            0  & |\nu/\bw| > 1/2 \\
            1/2& |\nu/\bw| = 1/2 \\
            1  & |\nu/\bw| < 1/2 \end{array} \right.$  \\ [6mm]
%
\hline
\multicolumn{2}{c}{Shift Theorem} \\
\hline \\ [-3mm]
%
$x'(t)=x(t + \tau)$ & 
$X'(\nu) = X(\nu) \exp(\Ci \pi \nu \tau)$  \\ [2mm]
%
\hline
\end{tabular}
\end{center}
\end{table}

%%%%%%%%%%%%%%%%%%%%%%%%%%%%%%%%%%%%%%%%%%%%%%%%%%%%%%%%%%%%%%%%%%%%%%%%%%%%%%%%%%%%%%%%

%%%%%%%%%%%%%%%%%%%%%%%%%%%%%%%%%%%%%%%%%%%%%%%%%%%%%%%%%%%%%%%%%%%%%%%%%%%%%%%%%%%%%%%%

\subsection{Narrow-band Approximation}
\label{app:narrow_band_approximation}

Presented with a scalar input signal, $e(t)$, the output of a one-dimensional linear time-invariant system is given by the convolution,
%
\begin{equation}
e'(t)=(j*e)(t) \equiv \int_{-\infty}^{\infty}{x(\tau)j(t-\tau)d\tau}.
\label{eqn:convolution}
\end{equation}
%
where $j(t)$ is the one-dimensional impulse response function.  In a two-dimensional
system, each output signal is given by a linear combination of the input
signals,
\begin{equation}
\begin{split}
e_1'(t)=(j_{00}*e_0)(t) + (j_{01}*e_1)(t), \\
e_2'(t)=(j_{10}*e_0)(t) + (j_{11}*e_1)(t).
\end{split}
\label{eqn:convolution2}
\end{equation}

Let $e_0(t)$ and $e_1(t)$ be the components of the electric field vector $\Jcol{e}(t)$ and $j_{mn}(t)$ be the elements of the $2\times2$ impulse response matrix, $\JM{j}(t)$.
%
By the convolution theorem, \eqn{convolution2} is equivalent to
\begin{equation}
\Jcol{E}'(\nu)=\JM{J}(\nu)\Jcol{E}(\nu),
\label{eqn:convolution_nu}
\end{equation}
where $\JM{J}(\nu)$ and $\Jcol{E}(\nu)$ are the Fourier transforms of $\JM{j}(t)$ and $\Jcol{e}(t)$. Under the assumption that $\JM{J}(\nu)$ is
approximately constant over a sufficiently narrow range of frequencies, convolution reduces to multiplication by a Jones matrix. 
